\subsubsection{Inversión de Hensel en la torre aritmética y realización por historias}
\label{sec:ch-capacidad-hensel}

Antes de levantar las ramas a historias enriquecidas, conviene recuperar el
argumento finito que determina la capacidad de la torre aritmética. Se
incorpora la prueba completa del transductor no filtrado de la matriz; su
dominio se distingue del árbol superviviente.

\begin{proposition}[Conjugación de la torre de Hensel no filtrada]
\label{prop:ch-hensel-bruto}
Sean \(p\) primo, \(d\geq1\) y \(A\in\operatorname{GL}_d(\mathbb Z)\).
Para vectores fila, tomando representantes de residuos en
\(\{0,\ldots,p-1\}^d\), se define
\[
 x_tA=u_t+px_{t+1},\qquad u_t=x_tA\bmod p.
\]
La aplicación
\[
 \Psi_T:(\mathbb Z/p^T\mathbb Z)^d\longrightarrow
             (\mathbb F_p^d)^T,\qquad
 x_0\longmapsto(u_0,\ldots,u_{T-1})
\]
es biyectiva, con inversa
\[
 x_0\equiv\sum_{t=0}^{T-1}p^tu_tA^{-(t+1)}\pmod{p^T}.
\]
Las biyecciones conmutan con las reducciones y dan un homeomorfismo
\(\Psi:\mathbb Z_p^d\to(\mathbb F_p^d)^{\mathbb N}\), donde
\(\mathbb Z_p:=\varprojlim_T\mathbb Z/p^T\mathbb Z\). La actualización
\(x_t\mapsto x_{t+1}\) corresponde a suprimir el primer bloque.
\end{proposition}

\begin{proof}
La inversa entera de \(A\) existe porque su determinante es \(\pm1\).
Iterando \(x_t=(u_t+px_{t+1})A^{-1}\), se obtiene
\[
 x_0=\sum_{t=0}^{T-1}p^tu_tA^{-(t+1)}+p^Tx_TA^{-T}.
\]
Módulo \(p^T\), una palabra determina una única clase inicial. Dada una
palabra arbitraria, fijar \(x_T=0\) y reconstruir hacia atrás produce una
clase que la emite, lo que prueba la sobreyectividad. La misma fórmula
prueba compatibilidad al truncar. El paso al límite inverso da el
homeomorfismo y la relación con el desplazamiento de la cinta.
\end{proof}

Para \(p=3\) y \(d=6\), el alfabeto de bloques tiene \(3^6=729\)
elementos. Esta capacidad aritmética no se identifica por definición con
la realizabilidad de cada bloque tras cada prefijo superviviente. El
enlace con las aristas \(\operatorname{Ext}\) es precisamente el cometido
del lema~\ref{lem:elevacion-catalogal-nonadica}. Se mantienen así las dos
operaciones que distingue la matriz: representación en la torre bruta y
realización bajo las compatibilidades APP--TRIT--TPK.

